\chapter{A worked audit and the next research boundary}\label{ch:next}
We end by reviewing a complete illustrative record. The point is not to
execute the future engine on paper. It is to identify exactly which facts
a reader would need in order to check one event, and which conclusions
remain unavailable even after the arithmetic is correct.

\section{The declaration}
Use two cells, reference $(10,10)$ and scales $(2,2)$. Begin at reference
with zero capacity and zero external audit. A declared conservative forcing
increment $(4,-4)$ creates baseline $(14,6)$. Its potential is four.
Two accepted child actions each transfer one unit from the first cell to
the second. The aggregate increment is $(-2,2)$ and endpoint $(12,8)$.

This is an illustrative declaration, not a registered world. It does not
select a random law, timestep, horizon or performance target. The children
are fixed mathematical objects for the audit.

\section{The physical and field checks}
Total stock is twenty in every listed state and every stock is nonnegative.
The accepted source withdrawal is two, within its fourteen-unit baseline.
If there were another declared physical constraint, it would also need
checking; the absence of such a field from this example is not permission
to ignore it in another model.

At baseline the marginals are $(1,-1)$. Each child has increment $(-1,1)$,
and the aggregate curvature term is
$\tfrac12(-2,2)^TH(-2,2)=1$. The tangent group benefit is four, so
the exact group value is three. Endpoint evaluation agrees: $4-1=3$.

For the straight common path, each child attribution is $3/2$. Two
standalone singleton values would instead be $7/4$ each and sum to
$7/2$, overcounting the group by one-half. This discrepancy is not a
floating-point artifact; it is the exact shared-field cross term.

\section{The account check and its missing authority}
Suppose, conditionally, the adopted policy assigns both children to one
owner. That owner's capacity change would be three. If each child belongs
to a separate owner, each would receive $3/2$ under the same declared
path-to-capacity policy. Both pass zero-genesis affordability in this example.
Neither assignment is derived from the sum identity; the owner rule is a
separate input that the actual programme still must adopt.

The external audit is four. With total capacity three and endpoint potential
one, $V+\sum B-J=1+3-4=0$. This verifies the conditional aggregate
account. It does not verify the provenance of the owner declaration, the
sensor readings or the random selection that supposedly chose the group.

\section{A list of claims the record cannot support}
It cannot show that the action was sampled uniformly, because no random
stream or menu record has been supplied. It cannot show that the process
will keep recovering, because only one fixed event is described. It cannot
show fairness, because arithmetic closure does not define fairness. It
cannot show causal entitlement to the child amounts, because the path
decomposition is not an intervention study.

It also cannot establish that a runnable Gaussian engine exists. A future
implementation must faithfully create, execute and preserve these objects
under its accepted profile, not merely print the same numbers for one example.

\section{The claim ledger}
\begin{center}
\begin{tabular}{P{.41\linewidth}P{.46\linewidth}}
\toprule Claim & Current status in this edition\\\midrule
Finite Gaussian expansion & Derived exactly under fixed quadratic geometry.\\
Local affected-factor cancellation & General conditional identity.\\
Common-path sum & Exact decomposition; no ownership conclusion.\\
Onsager-type dissipation & Conditional on the specified flux law and domain.\\
Capacity invariant & Conditional on the proposed signed owner/update policy.\\
Gaussian runtime acceptance & Not established by this book.\\
Gaussian scientific outcomes & No new study executed.\\
Historical Gate 1D-C completion & Recorded by its separate committed manifest.\\
\bottomrule
\end{tabular}
\end{center}

\section{The next separately authorized stages}
First resolve the load-bearing event, owner/path, menu/locality, forcing
coupling and numerical-policy contracts. Then implement reviewed generic
reuse and isolated Gaussian evaluation. Validate pure foundation properties
before separately authorized transition fixtures. Only after that should a
scientific packet freeze the study configuration, comparators and
interpretation rules. Execution requires packet-specific permission.

This sequence does not require the author to decide every routine coding
detail. It identifies the decisions that change the scientific object.
Routine implementation can proceed within an accepted contract; it must
stop where a missing choice would create a different process.

Stage A and Stage B are separate study questions, not automatic continuation
flags. A single-shock outcome, including nonrecovery, needs a disposition
before a continuously forced regime is treated as the next approved study.
Neither stage is automatically SD-01, and no AWS launch or cost is authorized
by a manuscript explaining the programme.

\section{What this revision accomplished}
The first two parts now have a consistent active Gaussian vocabulary and a
visible connection to the original force--mobility--flux mathematics. The
historical record remains identifiable rather than being rewritten as new
evidence. The derivations and worked arithmetic can be reviewed now while
implementation and empirical claims remain at their actual status.

That is not a delay in doing science. It is the distinction that lets the
next experiment answer a clear question. A model can be implemented, tested
and found wanting without destroying the exact identities it was built to
respect. Conversely, correct identities do not guarantee the model deserves
institutional adoption. Both possibilities must remain open.
